\chapter{Results}
\label{ch:results}

\section{Early matched-budget behaviour}
\label{sec:early-results}

Table~\ref{tab:12h} reports the matched $12$\,h pair. Both models are far from
the accuracies reported in the literature, which is expected at this budget
(\S\ref{ch:design}); the comparison of interest is between the columns.

\begin{table}[htbp]
\centering
\caption{Matched $12$\,h comparison on PoseBusters ($209$ complexes,
$10$ sampling seeds per complex). ``Single seed'' averages the per-seed rate;
``best of $10$'' is the oracle over seeds and is an upper bound on what any
ranking function could achieve.}
\label{tab:12h}
\small
\begin{tabular}{@{}lcc@{}}
\toprule
 & SigmaFlow & SigmaDock \\
\midrule
RMSD $<2$\,\AA, single seed      & $4.4\%$  & $\mathbf{9.8\%}$ \\
RMSD $<2$\,\AA, best of $10$     & $29.2\%$ & $\mathbf{45.9\%}$ \\
Median RMSD, single seed         & $4.78$\,\AA & $\mathbf{4.48}$\,\AA \\
Median RMSD, best of $10$        & $2.62$\,\AA & $\mathbf{2.10}$\,\AA \\
\bottomrule
\end{tabular}
\end{table}

SigmaDock leads on every entry. This is the honest starting point of the
empirical account, and it is stated before any interpretation.

Two features of the table deserve comment. First, the gap between the
single-seed and best-of-$10$ rows is very large for both arms: for SigmaFlow,
$4.4\%$ against $29.2\%$. Most of the achievable accuracy at this budget is
present in the sample set and lost at selection --- a fact that motivates
\S\ref{sec:confidence-results}. Second, the two arms differ by roughly a
factor of two throughout, rather than by a constant offset, which suggests a
difference in kind rather than in degree. \S\ref{sec:trans-vs-rot} localises
it.

\section{Training trajectory at matched compute}
\label{sec:traj-results}

\begin{reserved}{Training trajectory, $72$\,h matched pair}
\textbf{Experiment.} Two runs, SigmaFlow and SigmaDock, identical schedule
annealed to the full budget, checkpointed at $6,12,18,24,36,48,60,72$\,h.
Cheap evaluation on a fixed $60$-complex subset at every snapshot; full
$209$-complex, $3$-seed evaluation at $24$, $48$ and $72$\,h.

\textbf{Expected output.} Figure~\ref{fig:trajectory} (four panels) and a
table of annealed endpoints with bootstrap intervals.

\textbf{Metrics.} Median fragment rotation error against the Haar reference;
enrichment below $\Ang{30}$; median centroid displacement; median RMSD raw and
aligned; rotational and translational validation loss separately.

\textbf{Question answered.} Prediction~P3 of \S\ref{sec:predictions}: does the
rotational deficit narrow, persist, or vanish with additional training?

\textbf{Interpretation rules, fixed in advance.}
(i) If SigmaFlow's rotational enrichment rises towards SigmaDock's, P3 is
refuted and the correct reading is \emph{slower convergence}; the mechanism of
Chapter~\ref{ch:asymmetry} then explains the rate rather than a limit, and
\S\ref{sec:verdict} says so.
(ii) If SigmaFlow's enrichment stays flat at the Haar level while SigmaDock's
rises, P3 is supported and the deficit is attributed to the construction.
(iii) If both stay near Haar, no claim about the mechanism is licensed at this
budget and the honest reading is that neither model has learned orientation;
\S\ref{sec:limits} must then lead with that.
Reserved: \(\approx 700\) words.
\end{reserved}

\figplaceholder{fig:trajectory}
{Training trajectory at matched compute. Curves are snapshots from a single
annealed schedule; isolated markers are independently annealed endpoints
(\S\ref{sec:snapshot-caveat}).}
{Four panels sharing an $x$-axis of optimisation step (secondary axis:
wall-clock hours). Both arms on every panel, distinguished by colour.
\textbf{(a)} median fragment rotation error, with the Haar reference
$\Ang{132.3}$ as a horizontal dashed line;
\textbf{(b)} fraction of fragments with rotation error below $\Ang{30}$, with
the Haar rate ($\approx 2\%$) as a dashed line;
\textbf{(c)} median fragment centroid displacement;
\textbf{(d)} median RMSD, raw and after superposition.
Bands: bootstrap intervals over complexes. Annealed endpoints ($12$\,h pair,
$72$\,h endpoint) as filled markers \emph{not} joined to the curves.
Panel (b) is the decisive one: it shows whether SigmaFlow ever separates from
the null.}

\section{Translation versus rotation}
\label{sec:trans-vs-rot}

The factor-wise decomposition of the $12$\,h models localises the gap
unambiguously. Table~\ref{tab:rotation} reports per-fragment rotation error
against the Haar null of \S\ref{sec:metrics}.

\begin{table}[htbp]
\centering
\caption{Per-fragment rotational accuracy at $12$\,h ($711$ fragments over
$209$ complexes, seed $0$), against a uniformly random orientation.}
\label{tab:rotation}
\small
\begin{tabular}{@{}lccc@{}}
\toprule
 & SigmaFlow & SigmaDock & Haar (random) \\
\midrule
Median rotation error       & $\Ang{138.2}$ & $\mathbf{\Ang{123.8}}$ & $\Ang{132.3}$ \\
Fraction below $\Ang{30}$   & $2.1\%$ & $\mathbf{5.8\%}$ & ${\approx}\,2\%$ \\
Fraction below $\Ang{60}$   & $8.4\%$ & $\mathbf{15.8\%}$ & ${\approx}\,10\%$ \\
\bottomrule
\end{tabular}
\end{table}

SigmaFlow's rotational channel is not merely weaker than SigmaDock's: it is
not distinguishable from guessing. Its enrichment below $\Ang{30}$ matches the
random rate, whereas SigmaDock achieves roughly a threefold enrichment at
identical architecture, data and budget.

The translational channel behaves quite differently. The fraction of fragments
placed within $2$\,\AA{} is $22.6\%$ for SigmaFlow against $29.3\%$ for
SigmaDock, and within $1$\,\AA{} $8.2\%$ against $9.2\%$ --- a modest deficit,
but signal in both cases rather than its absence. The difference between the
two factors is therefore qualitative, not a matter of degree, which is what
Chapter~\ref{ch:asymmetry} sets out to explain.

\figplaceholder{fig:rotdist}
{Distribution of per-fragment rotation error at $12$\,h, against the
random-orientation null.}
{Three overlaid kernel density estimates on $[0,\Ang{180}]$: SigmaFlow,
SigmaDock, and the analytic Haar density. Vertical lines at the three medians.
Inset: zoom on $[0,\Ang{60}]$ where the enrichment differences live.
The point of the figure is immediate legibility --- the SigmaFlow density
should be seen to sit on top of the Haar density while SigmaDock's is shifted
left. Replaces roughly 300 words. Requires the per-fragment error arrays,
which exist for the $12$\,h models; regenerate for the $72$\,h models.}

\begin{reserved}{Mechanism diagnostics over training}
\textbf{Experiment.} Evaluate predictions P1 and P2 of \S\ref{sec:predictions}
at each checkpoint of the long runs: the ratio
$\|\pred^{\mathrm{rot}}\|/\|\vfield_t^{\mathrm{rot}}\|$ and the cosine between
predicted and target rotational fields, both stratified by $t$.

\textbf{Prior evidence.} On the $12$\,h model the magnitude ratio is
approximately $0.47$ and the cosine at $t<0.1$ approximately $0.025$, both
consistent with the shrinkage argument of \S\ref{sec:shrinkage}.

\textbf{Question answered.} Whether shrinkage is an early-training artefact
that resolves, or a stable property of the optimum.

\textbf{Interpretation.} A magnitude ratio rising towards $1$ with training
refutes P1 and points to under-training; a ratio stable well below $1$
supports the identifiability reading. Reserved: \(\approx 300\) words within
\S\ref{sec:trans-vs-rot} (total for the section: \(\approx 600\)).
\end{reserved}

\section{Integration efficiency}
\label{sec:nfe}

The conditional paths of \S\ref{sec:sigmaflow-objective} are straight in the
translational factor and geodesic in the rotational one, so the marginal field
they induce should be nearly straight and few integration steps should
suffice. Table~\ref{tab:nfe} tests this over a forty-fold range.

\begin{table}[htbp]
\centering
\caption{SigmaFlow accuracy against the number of Euler steps
($209$ complexes, $3$ seeds). The inherited default is $25$.}
\label{tab:nfe}
\small
\begin{tabular}{@{}lcccc@{}}
\toprule
Steps & Median RMSD ($\pm$SD over seeds) & $<2$\,\AA & $<5$\,\AA & Aligned RMSD \\
\midrule
$5$   & $4.61 \pm 0.06$ & $3.7\%$ & $54.9\%$ & $2.11$ \\
$10$  & $4.78 \pm 0.22$ & $3.5\%$ & $53.0\%$ & $1.97$ \\
$15$  & $4.88 \pm 0.11$ & $4.0\%$ & $51.8\%$ & $1.95$ \\
$25$  & $4.76 \pm 0.19$ & $4.1\%$ & $53.3\%$ & $2.01$ \\
$50$  & $5.08 \pm 0.05$ & $4.0\%$ & $49.1\%$ & $1.91$ \\
$100$ & $5.17 \pm 0.18$ & $4.8\%$ & $48.8\%$ & $1.99$ \\
$200$ & $4.93 \pm 0.24$ & $3.7\%$ & $50.7\%$ & $1.99$ \\
\bottomrule
\end{tabular}
\end{table}

The curve is flat. Paired per complex against the default of $25$ steps, the
largest effect is at $5$ steps ($-0.199$\,\AA, $95\%$ CI $[-0.391, -0.007]$);
the upper limit is essentially at zero, and with six comparisons against a
common reference this does not survive any correction for multiplicity. The
defensible statement is therefore that \emph{no step count between $5$ and
$200$ is demonstrably better than any other}. For scale, a single complex
varies across three sampling seeds with a standard deviation of $1.36$\,\AA,
whereas its seed-averaged range across all seven step counts is $1.66$\,\AA{}
--- less than would be expected from pure noise over seven draws.

Two readings follow, and both are needed. Positively, SigmaFlow reaches at $5$
steps what it reaches at $25$, so sampling can be made five times cheaper at no
measurable cost --- the signature predicted for near-straight conditional
paths, and the clearest advantage of the substitution observed so far.
Negatively, integration resolution is thereby \emph{excluded} as an
explanation for the accuracy gap: the model integrates its field faithfully;
the field points in the wrong direction.

The claim licensed by this table is ``SigmaFlow does not need the inherited
$25$ steps'', not ``SigmaFlow needs fewer steps than SigmaDock''. The latter
requires the matching sweep on the diffusion arm.

\begin{reserved}{Diffusion-arm step sweep}
\textbf{Experiment.} The same seven step counts and three seeds on the
SigmaDock arm, run strictly serially so that wall-clock is interpretable.
\textbf{Expected output.} A second series in Figure~\ref{fig:nfe}, and a
two-sided efficiency claim. \textbf{Question answered.} Whether the
few-step property is specific to flow matching or shared.
\textbf{Interpretation.} If SigmaDock also tolerates $5$ steps, the efficiency
argument for flow matching collapses and must be withdrawn; if it degrades,
the advantage is established. Reserved: \(\approx 150\) words.
\end{reserved}

\figplaceholder{fig:nfe}
{Accuracy against number of function evaluations, both arms.}
{Median RMSD (left axis) against integration steps on a log $x$-axis, $5$ to
$200$. Bands from three seeds. Two series once the diffusion sweep exists.
A second panel may show wall-clock per pose against steps, but \emph{only}
from serially executed jobs --- the existing array was run with five
concurrent tasks and its timings are contaminated by resource contention.}

\section{Diagnostics and ablations}
\label{sec:diagnostics}

\subsection{The frame correction and its signature}

The defect of \S\ref{sec:frame-defect} was identified analytically, but its
behavioural signature confirms that the evaluation is sensitive to a known
real intervention. We measure a \emph{locality gap}: the difference between
the rotation error of a fragment and that of its neighbours, negative when a
model places neighbouring fragments more coherently than chance.

\begin{table}[htbp]
\centering
\caption{Locality gap before and after the frame correction
($n \approx 133$ complexes, bootstrap intervals over complexes).}
\label{tab:framefix}
\small
\begin{tabular}{@{}lccc@{}}
\toprule
Run & Locality gap & $95\%$ CI & Fraction $<0$ \\
\midrule
SigmaFlow $6$\,h, uncorrected & $+\Ang{1.1}$  & $[-4.2, +6.6]$   & $53\%$ \\
SigmaFlow $6$\,h, corrected   & $-\Ang{12.8}$ & $[-18.3, -7.2]$  & $68\%$ \\
SigmaFlow $12$\,h, corrected  & $-\Ang{15.0}$ & $[-21.6, -8.1]$  & $69\%$ \\
SigmaDock $12$\,h             & $-\Ang{20.4}$ & $[-26.0, -14.5]$ & $71\%$ \\
\bottomrule
\end{tabular}
\end{table}

The uncorrected model is indistinguishable from no local coherence; the
correction moves it decisively. This licenses the use of the same metric
elsewhere.

\subsection{What does not explain the gap}
\label{sec:excluded}

Three candidate explanations were tested and none survives.

\begin{itemize}[leftmargin=*,itemsep=2pt]
\item \textbf{Integration resolution} --- excluded by \S\ref{sec:nfe}.
\item \textbf{Molecular symmetry.} Symmetry-equivalent atom assignments could
      in principle inflate apparent rotation errors. Measured directly, they
      do not account for the observed distribution.
\item \textbf{Loss weighting and objective variants.} Three single-change
      ablations off the same baseline --- a time-weighted flow loss, a
      rotational loss in data space rather than tangent space, and an
      anchor-atom distance term --- all returned a locality gap of
      $-\Ang{12.8}$, statistically indistinguishable from the unmodified
      baseline. Table~\ref{tab:ablations} summarises.
\end{itemize}

\begin{table}[htbp]
\centering
\caption{Single-change ablations, each branched from the same baseline.}
\label{tab:ablations}
\small
\begin{tabular}{@{}llc@{}}
\toprule
Variant & Single change & Locality gap \\
\midrule
Baseline (corrected)   & --- & $-\Ang{15.0}$ \\
(a) time weighting     & $w(t)\|\pred - \vfield\|^2$ & $-\Ang{12.8}$ \\
(b) rotational loss in data space & loss on $\Rot$ rather than $\Log \Rot$ & $-\Ang{12.8}$ \\
(c) anchor-atom distance & auxiliary distance term & $-\Ang{12.8}$ \\
\bottomrule
\end{tabular}
\end{table}

That three independent modifications of the objective produce no change is
itself evidence: it is consistent with a model whose rotational channel has no
signal to reweight. One caveat is recorded --- the reconstructed fragmentation
differs between runs for a substantial minority of complexes, so bond-level
comparisons across these variants are less reliable than the fragment-level
rotational statistics used here.

\section{Flow-specific source distribution}
\label{sec:source-results}

\S\ref{sec:source-theory} established that $p_0$ is free under flow matching
and fixed under diffusion, and \S\ref{sec:shrinkage} argued that the
uninformativeness of $\Haar$ is what starves the rotational channel.
Prediction~P4 follows: concentrating the rotational source should improve
rotation specifically.

Two constraints shape the construction. It must be \emph{sampleable} and
tractable, for which a wrapped Gaussian on $\SOthree$ about a centre
$\Rot_{\mathrm{c}}$ suffices,
\begin{equation}
  \Rot_0 \;=\; \Rot_{\mathrm{c}} \Exp(\skewm{\xi}),
  \qquad \xi \sim \mathcal{N}(0, \sigma^2 I_3),
  \label{eq:conc-source}
\end{equation}
with right-multiplication chosen to match the body-frame convention of
\S\ref{sec:trivialisation}. And it must not leak: $\Rot_{\mathrm{c}}$ may
depend only on the receptor and the ligand graph, never on the crystallographic
pose. This is not a formality --- an earlier variant in this project was found
to place conformer fragments at crystal centroids, which supplies the answer
through the source, and the constraint is therefore enforced explicitly and
tested by a control in which \eqref{eq:conc-source} reduces to $\Haar$ and
must reproduce the baseline bit-for-bit.

\begin{reserved}{Informative-source experiment}
\textbf{Gate.} A CPU-only audit measures whether receptor pocket geometry
carries orientational information at all, by comparing the rotational distance
under a candidate heuristic centre against the Haar reference of
$\Ang{132.3}$. \emph{Pre-registered rule:} if the median distance is not below
the Haar reference by more than its bootstrap interval, the experiment is not
run and this section reports the audit as a negative result in
\(\approx 350\) words.

\textbf{Experiment, if gated through.} A $6$\,h screen with
\eqref{eq:conc-source} against a bit-identical $6$\,h Haar control from the
same code path, the two differing only in $p_0$.

\textbf{Expected output.} Table~\ref{tab:source} completed;
Figure~\ref{fig:source}.

\textbf{Metrics.} Median rotation error and enrichment below $\Ang{30}$
(primary, since P4 predicts a rotation-specific effect); centroid displacement
and RMSD (secondary, expected approximately unchanged).

\textbf{Interpretation.} A rotation-specific improvement supports the
mechanism causally rather than correlationally and is the strongest available
evidence for Chapter~\ref{ch:asymmetry}. An improvement in \emph{both} factors
would instead suggest a generic optimisation benefit and would weaken the
specific claim. No improvement refutes P4.
Reserved: \(\approx 600\) words.
\end{reserved}

\begin{table}[htbp]
\centering
\caption{Effect of the source distribution on the rotational channel, at
matched budget.}
\label{tab:source}
\small
\begin{tabular}{@{}lccc@{}}
\toprule
Source $p_0$ on $\SOthree$ & Median rotation error & $<\Ang{30}$ & Median RMSD \\
\midrule
Haar (baseline)         & \tbd & \tbd & \tbd \\
Wrapped Gaussian, $\sigma$ = \tbd & \tbd & \tbd & \tbd \\
\midrule
Haar reference (random) & $\Ang{132.3}$ & ${\approx}\,2\%$ & --- \\
\bottomrule
\end{tabular}
\end{table}

\figplaceholder{fig:source}
{Source distributions on $\SOthree$ and their effect on the conditional
target.}
{Two panels. \textbf{(a)} Distribution of the initial geodesic distance
$\|\Log(\Rot_0^\top\Rot_1)\|$ under the Haar source and under
\eqref{eq:conc-source}, with the Haar mean $\Ang{126.5}$ marked. By
Proposition~\ref{prop:t-invariance} this is also the distribution of the
target magnitude at every $t$, which is the point of the panel.
\textbf{(b)} Rotation-error densities for the two trained variants, as in
Figure~\ref{fig:rotdist}, with the Haar null. Contingent on the gate.}

\section{Confidence and pose ranking}
\label{sec:confidence-results}

Table~\ref{tab:12h} shows that selection, not generation, is where most
achievable accuracy is lost at this budget: $4.4\%$ against a $29.2\%$ oracle.
A ranking function that recovered even a third of that difference would exceed
any architectural effect measured in this thesis.

\begin{reserved}{Confidence head}
\textbf{Experiment.} A small head on the $\ell=0$ (rotation-invariant)
channel of the existing backbone, trained in a second stage to predict
$\Prob(\mathrm{RMSD} < 2\,\text{\AA})$ for a generated pose. Two-stage rather
than joint, because during generative training the network sees only $x_t$,
which at $t=1$ is exactly the crystal pose; there are no poses of varying
quality to learn from.

\textbf{Expected output.} Table~\ref{tab:confidence}; Figure~\ref{fig:conf}.

\textbf{Metrics.} AUROC, AUPRC, Brier score, expected calibration error;
Spearman correlation against RMSD; and the operational quantity, top-$1$-of-$K$
accuracy against random-of-$K$ and the oracle.

\textbf{Question answered.} Whether the inherited backbone's invariant
features carry pose-quality information at all.

\textbf{Scope note.} The comparison against SigmaDock's published heuristic
ranker requires an external scoring binary that is not currently available in
the compute environment; if it remains unavailable, this section reports
calibration and top-$1$-of-$K$ only, and says so. \emph{This extension is not
flow-specific} --- an equivalent head could be attached to SigmaDock --- and
must be presented as a property of the system, not of flow matching.
Reserved: \(\approx 300\) words; drops to a Future Work paragraph if the head
is not trained.
\end{reserved}

\begin{table}[htbp]
\centering
\caption{Confidence-head calibration and ranking performance.}
\label{tab:confidence}
\small
\begin{tabular}{@{}lcccc@{}}
\toprule
 & AUROC & Brier & Spearman vs.\ RMSD & Top-$1$-of-$10$ $<2$\,\AA \\
\midrule
Random selection    & $0.50$ & --- & $0$ & \tbd \\
Confidence head     & \tbd & \tbd & \tbd & \tbd \\
Oracle (best of $10$) & --- & --- & --- & $29.2\%$ \\
\bottomrule
\end{tabular}
\end{table}

\figplaceholder{fig:conf}
{Confidence-head behaviour.}
{Three panels: (a) reliability diagram with the diagonal and per-bin counts;
(b) top-$k$ accuracy against $k$, for random selection, the confidence head,
and the oracle, showing how much of the oracle gap is recovered;
(c) mean predicted confidence against integration time $t$ along the sampling
trajectory, testing whether confidence becomes informative before $t=1$ and
could in principle support early termination. Contingent.}
